\noindent\textbf{Scope and editorial position.}
This chapter covers all 37 catalogue records, 53--89, from September 1965 to
March 1966, in the order of the handwriting. Its 298 mathematical
occurrences include repeated explanations and short interventions. The plan
in record 57, the year heading in record 78 and Kushnirenko's title-only
record 79 are preserved without invented mathematical content.
The supplied 52-page bundle begins with two carry-over pages,
\GSpage{P000088--P000089}, already represented by 1964-50.11--24;
Kohn's new-year talk starts on \GSpage{P000090}. Actual page boundaries
therefore replace the catalogue's overlapping attachment ranges.

\paragraph{The main mathematical lines.}
The most sustained autumn sequence develops the Lefschetz formula for
elliptic complexes: symbols, Hodge decomposition, positive-eigenspace
cancellation, regularized traces and the local fixed-point contribution.
The 1964 Harvard notes of Atiyah and Bott already contained this programme;
the later journal publication does not make Gelfand's exposition a new
1965 theorem \citep{atiyahbott1964notes,atiyahbott1967lefschetz}.
Kazhdan's parallel sequence supplies affine algebraic geometry, coherent
sheaves and derived functors. Gelfand's characteristic-form discussions
interact with Arnold's accounts of singularities, normal forms, canonical
transformations and ideal-fluid motion.

The spring records bring several more direct publication matches.
Gelfand--Kirillov's 1966 paper contains the Weyl-field and growth-dimension
results of record 80. The affine-plane transform constructions of record 82
match the 1967 Gelfand--Graev--Shapiro paper. Lidskii's all-zero zeta
continuation and Gindikin's multiple-time parabolic construction likewise
match precise 1967 results \citep{gelfandkirillov1966corps,gelfand1967planes,
lidskii1967zeros,gindikin1967multitime}. Agranovich's calculus had already
appeared in his 1965 survey \citep{agranovich1965elliptic}. These different
publication histories are kept separate from literal correctness and from
claims of historical priority.

\begin{center}
\begin{tabular}{@{}clrr@{}}
\toprule Class & Meaning & Occurrences & Share\\\midrule
A & Published soon after by the speaker/Gelfand & 22 & 7.4\%\\
B & Developed by others & 7 & 2.3\%\\
C & Partly developed & 1 & 0.3\%\\
D & Not developed & 0 & 0.0\%\\
E & Already known before the talk & 246 & 82.6\%\\
F & Not identifiable from the notes & 22 & 7.4\%\\
\bottomrule
\end{tabular}
\end{center}

\paragraph{What the critical reading changes.}
The apparatus separates a faithful transcription of a mathematical mistake
from a mistaken transcription. Examples of the former include the omitted
logarithm in growth dimension and the two extraneous factors in the top's
cyclic momenta. A direct counterexample shows why the torus intersection
bound needs transversality: a Hamiltonian graph on \(T^2\) can have three
intersections, so the unqualified \(2^2\) claim fails. An ellipsoid likewise
disproves the fixed-direction interpretation of the slice-volume formula.
An actual reading correction is decisive for Gindikin: the scan says
\ru{гиперб. оператор}, a hyperbolic operator, where the supplied text says
pseudo-operator. The original paper then provides an exact restricted
match. Restored hypotheses are labelled as reconstructions throughout.

\paragraph{Questions and a bounded development.}
The problem index follows 9 definite questions, including resolved ones.
There is one class-C occurrence: the request to find all automorphisms of
the Weyl division ring, GP-1965-08. Its unrestricted form is not
classified in the verified literature; the search did not locate a later
full solution. The development fixes \(D=Q(A_1(\mathbb C))\) and the
rational subfield \(\mathbb C(x)\). It gives complete proofs of its known
stabilizer normal form, an exact pole-and-residue criterion for inner
rational translations, the resulting translation subgroup in the outer
automorphism group, and finite-group cocycle removal. These are explicitly
restricted answers with no novelty claim. Two parameter runs and an
independent written review check the proofs and computations. The remaining
obstruction is moving the chosen rational subfield; no conclusion about
the whole automorphism group follows. There are no class-D items: the
unidentified fragments lack enough information to define research problems.

\paragraph{Continuities with 1964--65.}
The Weyl-rank question links directly to 1964-38.1, and the final
lattice-gas and thermodynamic-limit lines repeat 1964-3.1--2. The elliptic
sequence develops the language of 1964-28.2--5; the derived-functor
construction continues 1964-49.8. Smirnov's positive-definite-function
fragment is related to 1964-13.1, but its missing centrality and topology
prevent identification of the questions. These are links of different
strength, not assertions that every later talk settles an earlier problem.
The supplied previous index contains no problem identifiers, and its
cumulative bibliography is empty. New problem identifiers and checked
bibliography keys are therefore necessary. All present-status assessments
are dated 9 November 2026.
